\frontmatter \pagenumbering{Roman}  %% 采用大写罗马数字，缺省时为小写的

\chapter*{摘\texorpdfstring{\qquad}{} 要}\addcontentsline{toc}{chapter}{摘要}
\thispagestyle{main}

基于相位敏感光时域反射（phase-sensitive optical time-domain reflectometry，$\varphi$-OTDR）的分布式声学传感（distributed acoustic sensing，DAS）能够利用通信光纤连续测量沿线振动和动态应变，具有长距离、连续覆盖、抗电磁干扰和无源布设等特点。外差式$\varphi$-OTDR通过声光移频与相干接收保留瑞利后向散射光场信息，为扰动定位和波形恢复提供基础。然而，分立式光路的器件和连接点较多，系统体积、装调重复性及长期稳定性容易受到影响；脉冲发射、移频门控、偏振衰落、接收噪声和数字解调参数之间也存在耦合。

针对上述问题，本文围绕以集成SOA的ECL激光器为核心的外差式DAS开展研究。本文所称集成SOA的ECL激光器是一种在同一封装内集成外腔激光器（external cavity laser，ECL）与半导体光放大器（semiconductor optical amplifier，SOA）的ECL激光器，其中ECL产生窄线宽连续载波，SOA承担光功率放大和电脉冲形成。通过比较掺铒光纤放大器（erbium-doped fiber amplifier，EDFA）与SOA在增益、噪声、动态响应和集成方式方面的差异，说明本文选择SOA是面向无EDFA、快速电控和紧凑发射链路的适配性取舍。在此基础上，建立瑞利散射与距离映射、集成SOA的ECL激光器的增益和载流子--折射率耦合、声光调制器（acousto-optic modulator，AOM）动态响应、平衡相干探测及数字I/Q解调模型，研究集成SOA的ECL激光器驱动状态对输出脉冲、功率、消光比、放大自发辐射噪声和外差拍频的影响；同时设计SOA--AOM--DAQ协同门控时序，使有效光脉冲处于AOM稳定声光作用窗口内。

在硬件实现方面，将集成SOA的ECL激光器、AOM、$2\times2$光耦合器、环形器、偏振分束器（polarization beam splitter，PBS）和平衡光电探测器（balanced photodetector，BPD）进行模块级集成，完成具有脉冲发射、移频、收发分离、偏振分集和平衡探测功能的DAS一体化收发模组。本文以器件组成和功能拓扑相同的自搭建分立光路作为基线，在相同驱动、采集和数字处理流程下，从器件连接、外部接口和装调方式等方面与集成光路开展等条件比较，综合评价模块级集成的性能、工程代价与适用场景。

在传感性能验证方面，完成双偏振通道的幅值、偏置、I/Q方向和时延标定，采用统一的数字下变频、I/Q提取、复共轭差分、双偏振合成和滤波流程，并通过压电陶瓷振动实验评价扰动定位、频率和波形恢复能力。本文形成了“集成SOA的ECL激光器发射特性--收发光路集成--数字相位解调--DAS传感验证”的研究链条，可为外差式DAS系统的小型化和模块化实现提供参考。

\smallskip
\noindent {\bfseries 关键词：} 分布式声学传感；相位敏感光时域反射；集成SOA的ECL激光器；DAS一体化收发模组


\chapter*{\bfseries \sffamily Abstract}\addcontentsline{toc}{chapter}{Abstract}
\thispagestyle{main}

Distributed acoustic sensing (DAS) based on phase-sensitive optical time-domain reflectometry ($\varphi$-OTDR) enables continuous measurements of vibration and dynamic strain along standard telecommunication fibers. It provides long sensing range, continuous spatial coverage, immunity to electromagnetic interference, and passive deployment capability. In heterodyne $\varphi$-OTDR, an acousto-optic modulator introduces a frequency shift, while coherent reception preserves the amplitude and phase of the Rayleigh backscattered field for event localization and waveform reconstruction. However, a discrete optical setup generally contains numerous components and fiber connections, which may affect system size, alignment repeatability, and long-term stability. In addition, pulse generation, frequency-shift gating, polarization fading, receiver noise, and digital-demodulation parameters are mutually coupled, and the overall performance cannot be determined by a single device specification.

To address these issues, this thesis investigates a heterodyne DAS system employing an ECL laser with an integrated SOA. In this thesis, the device refers to an ECL laser that integrates an external cavity laser (ECL) and a semiconductor optical amplifier (SOA) within the same package. It provides two optical output ports: the LD port supplies continuous-wave reference light without SOA amplification, whereas the SOA port supplies amplified pulsed signal light formed by current gating. A comparison of erbium-doped fiber amplifiers (EDFAs) and SOAs in terms of gain, noise, dynamic response, and integration supports the selection of the SOA as an engineering tradeoff for a compact, rapidly controlled transmitter without an EDFA, rather than as a universally superior optical amplifier. A unified model describes Rayleigh backscattering and distance mapping, the gain and carrier--index coupling of the ECL laser with integrated SOA, acousto-optic modulator (AOM) frequency shifting and transient response, balanced coherent detection, and digital I/Q demodulation. The influence of the device driving state on the output pulse, optical power, extinction ratio, amplified spontaneous emission noise, and heterodyne beat note is investigated. An SOA--AOM--DAQ coordinated gating sequence is also designed so that the effective optical pulse lies within the stable acousto-optic interaction window of the AOM.

A DAS integrated transceiver module is further developed by assembling the ECL laser with integrated SOA, AOM, $2\times2$ optical coupler, optical circulator, polarization beam splitter (PBS), and balanced photodetector (BPD) at the module level. The module retains the pulse-generation, frequency-shifting, transmit--receive separation, polarization-diversity, and balanced-detection functions required by a heterodyne DAS system. A self-built discrete optical setup with the same components and functional topology is used as the baseline. The discrete and integrated implementations are compared using the same driving, acquisition, and digital-processing procedures. The benefits, costs, and applicability of module-level integration are assessed using optical, sensing, stability, size, interface, and alignment metrics.

For sensing-performance validation, the amplitude, offset, I/Q orientation, and delay of the dual-polarization channels are calibrated. A unified processing chain comprising digital downconversion, I/Q extraction, complex-conjugate spatial differencing, dual-polarization combining, and filtering is used to recover vibration location, frequency, and waveform. The thesis thereby establishes a research chain from the transmitter characteristics of the ECL laser with integrated SOA, through transmit--receive optical integration and digital phase demodulation, to DAS sensing validation, providing a reference for the miniaturization and modular implementation of heterodyne DAS systems.

\bigskip
\noindent {\bfseries Key Words:} distributed acoustic sensing; phase-sensitive optical time-domain reflectometry; ECL laser with integrated SOA; DAS integrated transceiver module
